\subsection*{把响应拆成固定事实和当前进度}

取章首的独立 3000 分购买。$A_J$ 固定 $S$、输入版本、原输入、P1/C1/T1、平台范围和 \code{pay-C1}。$R_J$ 取得成功证据后固定 X1、O1、完整 \code{provider\_result}、\code{payment\_result} 和 EV1 原报文。这些固定关系组成以后恢复时要保持的事实。P1 的付款进度、O1 的履约进度和查询中的恢复状态，则描述当前工作到了哪里。

例如，支付已经成功，EV1 尚未消费时，查询给出 \code{paid} 与 \code{unfulfilled}。F501 提交并且可信回报保存以后，履约进度成为 \code{fulfillment\_accepted}，\code{payment\_result} 继续包含原 O1/X1 和 3000。第十四章完整响应中的 \code{acceptance} 返回原 $S$、P/C、商品和报价快照，T1 及恢复责任另由 $A_J$ 保存。\code{payment\_result} 返回原 P/C/O、成功结论、平台范围与 X1、币种和金额，原 \code{provider\_result} 与 EV1 报文另由 $R_J$ 保存。若为重试重新生成 \code{occurred\_at}、另分配 EV2，消费者便会面对另一份事件。$R_J$ 保存 EV1 的原版本、发生时间和完整载荷，把事件恢复也放入了原结果关系。

\CFifteenFigure{ch15-04-observation}{固定受理与付款结果连接当前进度，查询分别说明各自的事实对象}{fig:ch15-observation}

为检查这些关系，把系统状态记为
\[
\Sigma=(J,\Gamma,D,W,L,Q).
\]
J 保存权威受理与结果，$\Gamma$ 保存会话资格，D 保存归属及代次，W 是购买库投影，L 是独立平台账本，Q 保存消费与履约事实。每个分量按所属事务域提交，不变量约束这些分量之间的对应关系。

\subsection*{每种允许迁移怎样保持单一绑定}

先看 J 中的身份关系。空初态没有绑定。$A_J$ 在同一权威仲裁中为一个完整 $S$ 建立一份不可覆盖的原输入与实际 P/C/T。已有 $S$ 的同参请求恢复原记录，异参请求返回冲突，二者都保持原绑定。第一次受理之前报价无效的来件没有形成购买，后来有效来件按首次规则处理。已经保存的业务拒绝则是这项 $S$ 的原结果，后来余额或价格改变也从原记录恢复。

然后检查 C 到平台的关系。C1 形成以后，所有发送者从同一记录读取原范围、\code{pay-C1} 和原支付输入。工作者领取代次决定谁能够提交受保护的本地进度，平台身份决定多次付款到达怎样合流。资源保存已经接受的领取代次水位，在同一事务中比较来件令牌、推进水位并写入受保护数据。资源已经接受 W2 的第 8 代以后，W1 携第 7 代恢复写入便被拒绝。已经离开 W1 的付款请求由平台原身份规则接住。

资金证据进入 $R_J$ 时，合流入口核对原 C、账户、币种、金额和平台身份。第一次可信成功分配实际 O1 与 EV1，重复的同内容证据恢复这份决定。另一份同 C 的不相容终态进入核查，原成功保持。可信拒绝同时覆盖既往无成功、全部已受理及在途路径已经终结、未来首次受理已关闭，才成为固定拒绝结果。

本地物化从这些决定复制关系。它可以插入缺失的 P/C/T/O/EV，可以把与 $A_J$ 完全一致的未决投影推进为 $R_J$ 的确定投影。已有行的归属、原输入、原结果或事件内容不一致时，整批回滚。消费事务 Q 再共同保存 I、B、实际 F 与原接收结果。每个边界都有一条具体规则，把新的事实连接到此前已经取得的身份。

\CFifteenFigure{ch15-05-invariants}{每种允许迁移保持自己的身份绑定，并核对来自上一个决定的完整关联}{fig:ch15-invariants}

RIFL 把完成记录与业务变化共同持久保存、重试必须找到原记录列为实现重复请求消除的核心问题~\cite{rifl}。本书的本地事务直接落实这一关系。外部链路再为每个独立提交域保存自己的完成事实，跨域缝隙由持久待办、固定身份和证据合流继续推进。

\subsection*{金额关系随业务动作一起扩展}

第一部分的封闭本地模型只有支出，所以余额等于初始余额减去已提交扣款。现在允许退款和其他入账，就要把收入放回同一等式。设 $b_0$ 是共同初始余额，$\mathcal G$ 是退款以外已发生的入账集合，$\mathcal D$ 是已发生的购买支出集合，$\mathcal F$ 是已发生的退款入账集合，$a$ 表示各项正金额。同一资金账本中的每项入账按用途进入其中一个集合，账户实际余额为
\[
b=b_0+\sum_{g\in\mathcal G}a(g)
       +\sum_{f\in\mathcal F}a(f)
       -\sum_{d\in\mathcal D}a(d).
\]

这里按资金账本中的成功事实计数。查询或回调尚未把某笔退款带回本地时，平台实际余额可能已经增加，本地知识则继续未决。把本地已确认退款总额直接当作真实平台余额，会在响应丢失的历史里漏掉这笔入账。

对于原收款 X1，退款准入还需要另一项关系。设 $A(X1)=3000$ 是原成功收款额，$\mathcal R_{\rm succ,J}(X1)$ 是 J 已确认成功的退款集合，$\mathcal R_{\rm pend,J}(X1)$ 是已经受理而尚无终态证据的退款集合。准入维护
\[
\sum_{f\in\mathcal R_{\rm succ,J}(X1)}a(f)
+\sum_{u\in\mathcal R_{\rm pend,J}(X1)}a(u)
\le A(X1).
\]

成功退款与未决占用分开相加，每项已受理退款恰在其中一组，或已经获得足以释放额度的最终拒绝。未决退款可能在平台已经成功，只是本地尚未知晓。它继续占用额度，后面取得成功证据时从 $\mathcal R_{\rm pend,J}$ 移到 $\mathcal R_{\rm succ,J}$，总占用保持。

\subsection*{退款不变量怎样归纳}

空初态中 $R_{\rm succ}=R_{\rm pend}=0$。首次受理金额 $x$ 前，守卫要求 $R_{\rm succ}+R_{\rm pend}+x\le A$，提交后 $R_{\rm pend}$ 增 $x$，所以不等式保持。同 RF 同参恢复、异参冲突和未知观察都不改变 $R_{\rm succ}/R_{\rm pend}$。成功把 $x$ 从 $R_{\rm pend}$ 移到 $R_{\rm succ}$，总量保持。足够的最终拒绝把 $R_{\rm pend}$ 减 $x$，总量下降。不同来源分别执行同一归纳。

这项证明依赖身份不重用、原金额不可覆盖、同来源预占共同仲裁以及终态证据正确。若已拒绝 RF 以后还能在平台成功，释放额度这一步就失去前提。若恢复旧备份丢掉 RF1 的未决预占，RF2 又可使用同一份额度。退款的权威绑定与占用因此也写入跨本地回退保存的 J，本地退款投影按原 RF 重建。

\CFifteenFigure{ch15-12-refund-occupancy}{成功与未决占用共同受原收款额约束，退款投影回退后从 J 恢复原额度责任}{fig:ch15-refund-occupancy}


例如 RF1 已确认成功 1000，RF2 未决 1500，可新增的退款额度为 500。账本中的净收款由真实成功退款计算，受理入口中的可用额度由已确认成功加全部未决占用计算。两项数值分别服务于资金核对与新责任仲裁。

% 图注补明：图中使用另一条独立历史。
\CFifteenFigure{ch15-19-money}{另一独立历史中 RF1 已退还 1000 而本地尚未知晓，真实资金与未决占用分别计算}{fig:ch15-money}

\subsection*{进展具体卡在哪一步}

让 C1 的平台请求已成功、响应丢失，本地未决。固定身份与单一绑定可以保护安全，进展还需要恢复扫描器找到 T1、取得执行机会、成功查询原身份，并在 J 与购买库提交结果。扫描器永远不被调度时，平台 X1 可以长期存在，P1 也可以长期显示核查中。

再让 EV1 已经写入 Outbox，却没有发布。发布者最终可运行、消息通道最终接受、失败后的原事件持续重试，才会让消费者获得来件。消费者的数据库最终可提交，竞争事务最终结束，Q 才形成 F501。ACK 永远丢失可能使投递继续增加，但 I/B 让后续投递回到原 F。

客户端获知又多一段条件。Alice 持久保存原 $S$ 与原输入，并在可用通道恢复后继续查询或重试，才能最终取得原结果。她的设备永久离线时，服务端可以完成 C/T，结果保存在查询入口，获知则等待设备重新通信。

\CFifteenFigure{ch15-20-progress}{平台提交、服务端保存与客户端获知分别依赖后续机会，原事实在等待期间保留}{fig:ch15-progress}

新增责任速度也会影响进展。若每分钟受理的未知付款多于恢复者能够核对的数量，旧任务可能不断被新任务挤到后面。恢复队列安排公平扫描，给未决数量和占用金额设置容量，容量达到边界时暂停相应新受理，为已经接受的工作保留持续推进的空间。

\subsection*{证据永久丢失时保存什么结论}

取同一份旧购买备份 B0，构造两条历史。H0 中，C1 的请求从未到达平台。H1 中，平台已经完成 X1，再丢失本地与独立关联。恢复者收到同一 $S$，读到同一 B0，普通平台查询又都暂时查空。它当前取得的观察完全相同。

若恢复者给两条历史都建立 C2，H1 可能增加另一笔付款。若都返回成功，H0 又没有相应资金事实。停止自动发送、保存原来可能存在的工作并继续取证，是当前信息下可执行的共同动作。后续完整账本、可靠回执或平台终结证据一旦到来，两条历史才有不同的继续依据。

增加恢复线程只能更快地重复相同查询。部署保存哪份跨故障证据，决定了恢复者以后能否重新区分这些历史。独立 J 保存 $S\to$实际 C$\to$原平台身份，让恢复者能够从原请求重新取得查询与恢复的对象。

\CFifteenFigure{ch15-21-indistinguishable}{两条历史给恢复者相同观察，可信资金证据或旧路径终结证据才能区分后续动作}{fig:ch15-indistinguishable}
\FloatBarrier
